%%
%% Appendix B: Per-category tables + per-class complementarity
%%
% Source: thesis/tables/tab_percat_{ucf,xd}.tex (Class R, generated by
%   scripts/thesis_per_category_tables.py from tracked
%   results/{ucf_gated_fusion_giant,xd_gated_fusion_so400m}_s{42,123,2024}/
%   per_category.csv, post-lambda0/post-h1 — the ONLY per-category source).
% Per-class complementarity: results/_analysis_2026-06-10/
%   pri7_per_class_complementarity.{csv,md} (tracked).
\chapter{Per-Category Results and Per-Class Complementarity}
\label{app:b}

This appendix provides the complete per-category tables for the two headline
configurations summarized in Section~\ref{sec:percat}, followed by the
per-class complementarity analysis referenced from
Chapters~\ref{ch:05} and~\ref{ch:07}. All per-category values are computed
from the current canonical run directories (three seeds, the final
configuration of Chapter~\ref{ch:04}); no earlier-era per-category numbers
appear anywhere in this thesis.

\section{Full Per-Category Tables}
\label{sec:appb-tables}

% Source: generated by scripts/thesis_per_category_tables.py (CPU-only, tracked).
% Data: results/ucf_gated_fusion_giant_s{42,123,2024}/per_category.csv (tracked, post-lambda0/post-h1).
% 3-seed mean+/-std, frame-level AUC (%).
\begin{table}[t]
\centering
\caption{Per-category frame-level AUC (\%) on UCF-Crime for the headline Gated Fusion variant (SigLIP2 Giant backbone). Mean$\pm$std over three seeds (42, 123, 2024).}
\label{tab:percat-ucf}
\begin{tabular}{lr}
\toprule
Category & AUC (\%) \\
\midrule
Abuse & 96.05$\pm$0.96 \\
Arrest & 90.97$\pm$0.25 \\
Arson & 96.88$\pm$0.55 \\
Assault & 97.18$\pm$0.45 \\
Burglary & 94.84$\pm$0.37 \\
Explosion & 88.75$\pm$1.04 \\
Fighting & 93.92$\pm$0.37 \\
RoadAccidents & 87.14$\pm$0.58 \\
Robbery & 88.64$\pm$0.71 \\
Shooting & 95.67$\pm$0.33 \\
Shoplifting & 91.81$\pm$0.65 \\
Stealing & 99.34$\pm$0.13 \\
Vandalism & 97.47$\pm$0.38 \\
\bottomrule
\end{tabular}
\end{table}

% Source: generated by scripts/thesis_per_category_tables.py (CPU-only, tracked).
% Data: results/xd_gated_fusion_so400m_s{42,123,2024}/per_category.csv (tracked, post-lambda0/post-h1).
% 3-seed mean+/-std, frame-level AP (%).
\begin{table}[t]
\centering
\caption{Per-category frame-level AP (\%) on XD-Violence for the headline Gated Fusion variant (SigLIP2 SO400M backbone). Mean$\pm$std over three seeds (42, 123, 2024).}
\label{tab:percat-xd}
\begin{tabular}{llr}
\toprule
Code & Category & AP (\%) \\
\midrule
B1 & Fighting & 81.13$\pm$0.23 \\
B2 & Shooting & 56.84$\pm$0.95 \\
B4 & Riot & 94.03$\pm$0.27 \\
B5 & Abuse & 43.57$\pm$3.53 \\
B6 & Car Accident & 48.16$\pm$1.18 \\
G & Explosion & 72.22$\pm$1.11 \\
\bottomrule
\end{tabular}
\end{table}


Table~\ref{tab:percat-ucf} lists per-category frame-level AUC on UCF-Crime
for the headline Gated Fusion variant (SigLIP2 Giant backbone), and
Table~\ref{tab:percat-xd} lists per-category frame-level AP on XD-Violence
for the headline variant there (SigLIP2 SO400M backbone); both report
mean$\pm$std over seeds 42, 123, and 2024.

Three observations organize the UCF-Crime table. First, the spread across
categories is large---from Stealing (99.34$\pm$0.13\% AUC) down to
RoadAccidents (87.14$\pm$0.58\%)---and the weak tail (RoadAccidents,
Robbery, Explosion) is exactly the abrupt, short-event category set that the
failure analysis of Section~\ref{sec:failure} identifies independently from
score traces. Second, seed variance is small in most categories (std below
one point everywhere except Explosion at 1.04), so the per-category ordering
is stable across training runs. Third, several categories have very few
evaluated test videos (Section~\ref{sec:appb-pri7} tabulates the counts), so
the per-category values themselves carry sampling uncertainty well beyond
the seed-std column---across-category comparisons should be read
qualitatively, not as precise rankings.

On XD-Violence the picture is more polarized. Riot (94.03$\pm$0.27\% AP) and
Fighting (81.13$\pm$0.23\%) are strong; Abuse (43.57$\pm$3.53\%) and Car
Accident (48.16$\pm$1.18\%) are weak. Because AP is sensitive to
positive-frame prevalence, categories whose events are brief relative to
their videos (Abuse especially) score far lower in AP than their ranking
quality alone would suggest; the same prevalence effect appears in extreme
form in the RTFM reproduction's per-category profile
(Appendix~\ref{app:a}). Abuse is also the least seed-stable category
(std 3.53), consistent with its small sample.

\section{Per-Class Complementarity (Does the Skeleton Help Motion
Classes More?)}
\label{sec:appb-pri7}

The complementarity result of Chapter~\ref{ch:05}---gated fusion meets or
exceeds its visual-only counterpart in all eight dataset$\times$backbone
configurations, with a mean gain of $+0.6$ points---is an aggregate
statement. A natural mechanistic follow-up is whether the skeleton stream
helps most where its inductive bias says it should: on categories defined by
articulated human motion. This section tests that hypothesis at per-class
granularity on UCF-Crime.

\paragraph{Setup.} For the UCF-Crime headline backbone (SigLIP2 Giant), we
compute per-category frame-level AUC for Gated Fusion and for its
Visual-Only counterpart, average each over the three seeds, and take the
difference $\Delta$ (positive = skeleton helps that category). The 13
categories are grouped \emph{a priori} into a human-motion group (HM: events
defined by body dynamics, e.g.\ Fighting, Assault, Shooting) and an
appearance group (APP: events recognizable from scene appearance with little
articulated-motion signature: Arson, RoadAccidents, Explosion).

% Source: results/_analysis_2026-06-10/pri7_per_class_complementarity.csv
% (tracked; 3-seed means, giant backbone). n = evaluated anomaly test videos
% per category (the 254-video post-exclusion test set; anomaly n sums to 128).
\begin{table}[t]
\centering
\caption{Per-class skeleton gain on UCF-Crime (Gated Fusion minus
Visual-Only, SigLIP2 Giant, 3-seed mean), sorted by $\Delta$. HM =
human-motion group, APP = appearance group; $n$ = evaluated anomaly test
videos in the category; $\sigma$ = std of $\Delta$ across seeds.}
\label{tab:appb-pri7}
\begin{tabular}{llrrrrr}
\toprule
Category & Group & $n$ & Fused AUC (\%) & Visual AUC (\%) & $\Delta$ (pp) & $\sigma$ (pp) \\
\midrule
Fighting      & HM  & 5  & 93.92 & 92.20 & $+1.72$ & 0.67 \\
Arrest        & HM  & 5  & 90.97 & 89.25 & $+1.72$ & 1.07 \\
Shooting      & HM  & 22 & 95.67 & 94.62 & $+1.05$ & 0.20 \\
Burglary      & HM  & 12 & 94.84 & 94.11 & $+0.73$ & 0.72 \\
Shoplifting   & HM  & 19 & 91.81 & 91.26 & $+0.55$ & 0.40 \\
Vandalism     & HM  & 5  & 97.47 & 97.24 & $+0.22$ & 0.33 \\
Stealing      & HM  & 5  & 99.34 & 99.16 & $+0.18$ & 0.09 \\
Arson         & APP & 9  & 96.88 & 97.07 & $-0.19$ & 0.50 \\
RoadAccidents & APP & 16 & 87.14 & 87.35 & $-0.21$ & 2.42 \\
Explosion     & APP & 20 & 88.75 & 89.24 & $-0.48$ & 1.90 \\
Abuse         & HM  & 2  & 96.05 & 96.84 & $-0.80$ & 1.05 \\
Assault       & HM  & 3  & 97.18 & 98.05 & $-0.87$ & 0.43 \\
Robbery       & HM  & 5  & 88.64 & 90.10 & $-1.46$ & 1.97 \\
\bottomrule
\end{tabular}
\end{table}

\paragraph{Results.} Table~\ref{tab:appb-pri7} gives the full breakdown. At
the group level, the human-motion classes gain $+0.31\pm1.08$\,pp on average
(7/10 classes positive), while the appearance classes lose
$-0.29\pm0.16$\,pp (0/3 positive); the group-mean difference is
HM$-$APP $= +0.60$\,pp. The cleanest single class is Shooting at
$+1.05\pm0.20$\,pp---notable because it is also the best-sampled category
($n=22$), so its estimate is the least noisy in the table. The direction of
every one of these facts matches the mechanistic story: skeleton dynamics
help where events are defined by body motion, and contribute nothing (or
mild noise) where they are not.

\paragraph{Why we do not lean on this.} The evidence is
directionally supportive but mixed/noisy, and we treat it as descriptive
rather than confirmatory, for three reasons. First, the within-group spread dwarfs the
between-group gap: the HM group's class-to-class std (1.08\,pp) is nearly
twice the $+0.60$\,pp group difference. Second, the three largest negative
deltas---Robbery ($-1.46$), Assault ($-0.87$), Abuse ($-0.80$)---all sit in
the human-motion group, and all three are tiny-sample categories ($n=5$, 3,
and 2 evaluated anomaly videos respectively), where a single video can swing
the class AUC by several points; the $\sigma$ column captures only
seed-to-seed training noise, not this (larger) sampling noise. Third, the
group mean is unweighted, so a two-video class influences it as much as the
22-video Shooting class. We therefore report the per-class analysis for
transparency, note that its direction is consistent with the
complementarity mechanism argued in Chapter~\ref{ch:05}, and rest the
headline complementarity claim on the aggregate 8-configuration comparison,
not on this breakdown.
